\documentclass[11pt,a4paper]{article}
\usepackage{amsmath,amssymb,amsthm}
\usepackage[margin=2.5cm]{geometry}
\usepackage{hyperref}

\newtheorem{definition}{Definition}[section]
\newtheorem{theorem}[definition]{Theorem}
\newtheorem{proposition}[definition]{Proposition}
\newtheorem{corollary}[definition]{Corollary}
\newtheorem{remark}[definition]{Remark}
\newtheorem{conjecture}[definition]{Conjecture}

\title{Hyperseed Formalization:\\
  Reflections on the World Crystal}
\author{ProtomegaTron (automated formalization)}
\date{2026-09-19 (deepened v2)}

\begin{document}
\maketitle

\begin{abstract}
We formalize Ben Goertzel's ``Reflections on the World Crystal''
(2022-11-03) within the Hyperseed ontology, incorporating d-calculus
extensions from Hyperseed v2. The world crystal is modeled as a lattice
base space with discrete fiber, the emergence of continuity as fiber
coarse-graining via a surjective bundle morphism, and consciousness as
a fiber resonance pattern whose substrate independence follows from
fiber isomorphism under coarse-graining.
\end{abstract}

\tableofcontents

\section{Lattice base space}

\begin{definition}[Lattice base space]
The \emph{world crystal base space} $B_{\text{lattice}} = \{b_i\}_{i \in I}$
is a discrete set of fundamental events/nodes. Each site $b_i$ carries
a fiber $E_{b_i}$, forming a discrete fiber bundle
$\pi: E_{\text{lattice}} \to B_{\text{lattice}}$.
\end{definition}

\begin{definition}[Lattice adjacency]
The lattice structure is specified by an adjacency function
$g: I \times I \to \{0,1\}$ where $g_{ij} = 1$ if sites $i$ and $j$
are adjacent. This gives $B_{\text{lattice}}$ the structure of a graph.
\end{definition}

\section{Emergence via coarse-graining}

\begin{definition}[Coarse-graining morphism --- claims 5, 9, 10]
The \emph{coarse-graining morphism} is a surjective bundle map:
\[
  \text{CG}: (E_{\text{lattice}}, B_{\text{lattice}})
    \to (E_{\text{manifold}}, B_{\text{manifold}})
\]
where $B_{\text{manifold}}$ is the quotient of $B_{\text{lattice}}$ by a
scale-dependent equivalence relation $\sim_\ell$ that identifies sites
indistinguishable at scale $\ell$:
\[
  B_{\text{manifold}}(\ell) = B_{\text{lattice}} / \sim_\ell
\]
The continuum limit is $B_{\text{manifold}} = \lim_{\ell \to \infty}
B_{\text{lattice}} / \sim_\ell$.
\end{definition}

\begin{proposition}[Information loss --- claim 11]
\label{prop:info-loss}
Coarse-graining is lossy:
\[
  \dim(E_{\text{lattice}}) > \dim(E_{\text{manifold}})
\]
The information lost in coarse-graining corresponds to lattice-scale
degrees of freedom that manifest as quantum uncertainty at the continuum
level. The ``smooth'' continuum fiber is an average that suppresses
lattice-scale fluctuations.
\end{proposition}

\begin{proposition}[Emergent smoothness --- claims 1, 2, 3, 4, 5]
\label{prop:smoothness}
The coarse-grained bundle $(E_{\text{manifold}}, B_{\text{manifold}})$
is differentiable (admits smooth sections) even though the underlying
lattice bundle is discrete:
\[
  B_{\text{lattice}} \text{ discrete} \quad \xrightarrow{\text{CG}}
  \quad B_{\text{manifold}} \text{ smooth}
\]
This is the world crystal pattern: discrete atoms of geometry
$\to$ smooth spacetime, shared by Wolfram, LQG, and causal set approaches.
\end{proposition}

\section{Consciousness as fiber resonance}

\begin{definition}[Fiber resonance pattern --- claims 6, 12]
A \emph{consciousness resonance pattern} is a coherent section of the
fiber bundle over a connected sublattice $S \subset B_{\text{lattice}}$:
\[
  \psi: S \to E_{\text{lattice}}, \quad
  \text{coh}(\psi) = \frac{|\langle \psi(b_i), P_{ij} \psi(b_j) \rangle|}
    {\|\psi(b_i)\| \|\psi(b_j)\|} \geq \theta
  \quad \forall\, (i,j) \text{ adjacent in } S
\]
where $P_{ij}$ is parallel transport and $\theta$ is a coherence
threshold. High coherence = integrated, correlated fiber dynamics =
consciousness.
\end{definition}

\begin{proposition}[Substrate independence --- claims 7, 13, 18]
\label{prop:substrate}
Two lattices $B_1$ and $B_2$ support the same consciousness if their
coarse-grained fiber bundles are isomorphic on the resonance sublattice:
\[
  \text{CG}(E_1, S_1) \cong \text{CG}(E_2, S_2)
  \implies \text{same consciousness pattern}
\]
This is the fiber-geometric content of substrate independence: different
lattice substrates (biological neurons, silicon transistors, quantum
dots) produce isomorphic coarse-grained fiber bundles when they support
the same resonance pattern.
\end{proposition}

\section{d-calculus on the lattice (Hyperseed v2)}

\begin{definition}[Lattice covariant derivative --- claim 14]
The \emph{lattice covariant derivative} replaces the continuum
$\nabla$ with a finite difference operator:
\[
  (\nabla_{\text{lat}} \sigma)(b_i) = \sum_{j: g_{ij}=1}
    \left[\sigma(b_j) - P_{ij}\, \sigma(b_i)\right]
\]
where $P_{ij}: E_{b_i} \to E_{b_j}$ is discrete parallel transport
(a fiber isomorphism between adjacent sites). This is the fundamental
d-calculus object in the world crystal.
\end{definition}

\begin{proposition}[Plaquette curvature --- claim 15]
\label{prop:plaquette}
Curvature on the lattice is measured by holonomy around elementary
plaquettes --- the smallest closed loops $\gamma = (b_1, b_2, b_3, b_4)$
in the lattice:
\[
  F_{\gamma} = P_{12}\, P_{23}\, P_{34}\, P_{41} - \mathrm{id}_{E_{b_1}}
\]
Non-zero $F_\gamma$ indicates non-trivial geometry at the lattice scale.
In LQG, $F_\gamma$ is literally the curvature of spacetime at the
smallest resolvable scale.
\end{proposition}

\begin{proposition}[Continuum curvature emergence --- claim 16]
\label{prop:continuum-curv}
The continuum Riemann curvature emerges as the statistical average of
plaquette curvatures:
\[
  F_{\text{continuum}}(x) = \lim_{N \to \infty}
    \frac{1}{N} \sum_{\gamma \ni x} F_\gamma
\]
where the sum is over the $N$ plaquettes containing the point $x$ in the
coarse-grained manifold. Smooth curvature = statistical average of
discrete curvature.
\end{proposition}

\begin{proposition}[Consciousness curvature --- claim 17]
\label{prop:consciousness-curv}
The \emph{consciousness curvature} is the fiber curvature restricted
to the consciousness resonance sublattice $S$:
\[
  F_{\text{conscious}}(b_i) = F_\nabla(b_i)\big|_{S}
  \quad \text{for } b_i \in S
\]
Non-zero consciousness curvature means that the conscious experience
varies non-trivially across the substrate --- different regions of $S$
contribute differently to the integrated experience. This is the
geometric content of ``qualia vary by neural substrate.''
\end{proposition}

\begin{conjecture}[Computational equivalence as fiber isomorphism --- claim 18]
Wolfram's principle of computational equivalence, in fiber-geometric
language, states:
\[
  \forall\, B_1, B_2 \text{ (sufficiently complex lattices)}: \quad
  \text{CG}(E_1, B_1) \cong \text{CG}(E_2, B_2)
\]
All sufficiently complex lattice dynamics produce isomorphic fiber
bundles at the coarse-grained level. This is why different computational
substrates can support the same intelligence and consciousness ---
their lattice details differ, but their emergent fiber geometry converges.
\end{conjecture}

\end{document}
